\section{Источники погрешности в математическом моделировании}
\label{sec:error-sources-in-modeling}

Результат вычислительного моделирования является приближённым по нескольким
принципиально различным причинам. Удобно разделять ошибку по этапам
перехода
\[
\text{реальный объект}
\to
\text{математическая модель}
\to
\text{численный метод}
\to
\text{вычисления в ЭВМ}.
\]

Пусть $y_{\rm real}$ --- интересующая характеристика реального объекта,
$y$ --- точное решение принятой математической модели, $y_h$ --- точный
результат дискретного метода при параметре $h$, а $\widehat y_h$ --- реально
вычисленное в конечной арифметике значение. Тогда формально
\[
\widehat y_h-y_{\rm real}
=
\underbrace{(y-y_{\rm real})}_{\text{погрешность модели и данных}}
+
\underbrace{(y_h-y)}_{\text{погрешность метода}}
+
\underbrace{(\widehat y_h-y_h)}_{\text{вычислительная погрешность}}.
\]
Это разложение подчёркивает, что уменьшение одной составляющей не устраняет
остальные.

\subsection{Неустранимая погрешность и погрешность математической модели}

Источниками первой группы являются
\begin{itemize}
    \item идеализация реального объекта и отброшенные физические эффекты;
    \item неточно известные параметры, коэффициенты, начальные и граничные
    данные;
    \item погрешности измерений;
    \item неполное знание внешних воздействий.
\end{itemize}

Погрешность, связанную с неточностью исходной физико-математической
постановки, часто называют \textbf{неустранимой на этапе численного
решения}: её нельзя уменьшить простым сгущением сетки. Для этого нужна
более адекватная модель или более точные данные.

\subsection{Погрешность численного метода}

Точный оператор задачи заменяется приближённым. Если
\[
\|y_h-y\|\le Ch^p,
\]
то говорят о сходимости порядка $p$ относительно параметра $h$.
Для методов решения дифференциальных уравнений необходимо различать
\textbf{аппроксимацию}, \textbf{устойчивость} и \textbf{сходимость}.
Малый локальный дефект схемы сам по себе не гарантирует малую глобальную
ошибку: возмущения не должны неограниченно усиливаться алгоритмом.

Практически погрешность дискретизации оценивают, например, по серии расчётов
$h$, $h/2$, $h/4$ или с помощью экстраполяции Ричардсона.

\subsection{Абсолютная и относительная погрешности}

Если $a$ --- точное значение, а $\widehat a$ --- приближённое, то
\[
\Delta a=|\widehat a-a|
\]
называется абсолютной погрешностью, а при $a\ne0$
\[
\delta a=\frac{|\widehat a-a|}{|a|}
\]
--- относительной погрешностью.

Для функции $y=f(a_1,\ldots,a_n)$ при малых возмущениях входных данных
используется линеаризация
\[
\Delta y\approx
\sum_{i=1}^n
\frac{\partial f}{\partial a_i}\Delta a_i.
\]
Она показывает, какие параметры сильнее всего влияют на результат, и
лежит в основе анализа чувствительности.

\subsection{Обусловленность задачи и устойчивость алгоритма}

Важно различать свойства \textbf{задачи} и \textbf{алгоритма}.
Обусловленность характеризует чувствительность точного решения к малым
изменениям исходных данных. Для СЛАУ
\[
Ax=b
\]
в согласованной норме используется число обусловленности
\[
\boxed{\kappa(A)=\|A\|\,\|A^{-1}\|}.
\]
Большое $\kappa(A)$ означает, что даже точный алгоритм не может устранить
сильное влияние ошибок данных.

Численно устойчивый алгоритм, напротив, не должен существенно усиливать
малые погрешности, возникающие в ходе вычислений. Поэтому плохо
обусловленная задача и неустойчивый алгоритм --- разные причины потери
точности.

\subsection{Вычисления с плавающей точкой}

В ЭВМ вещественные числа хранятся с конечным числом значащих разрядов.
В нормализованной двоичной системе число имеет вид
\[
x=\pm m\,2^e,
\]
где мантисса $m$ и порядок $e$ принадлежат конечным множествам.
Следовательно, большинство вещественных чисел приходится округлять.

Для основных арифметических операций при отсутствии переполнения и
потери порядка стандартная модель округления имеет вид
\[
\boxed{
\operatorname{fl}(x\circ y)
=(x\circ y)(1+\delta),
\qquad |\delta|\le u,
}
\]
где $u$ --- единичная погрешность округления, а $\circ$ обозначает одну из
основных арифметических операций. Машинный эпсилон
$\varepsilon_{\rm mach}$ обычно определяют как расстояние от $1$ до
следующего представимого числа; при округлении к ближайшему $u$ имеет тот
же порядок и обычно равен половине этого расстояния.

\subsection{Накопление и потеря значащих цифр}

Малые ошибки отдельных операций могут накапливаться. Особенно опасно
\textbf{вычитание близких чисел}: ведущие значащие цифры сокращаются, и
относительная ошибка резко возрастает. Например, выражение
\[
\sqrt{x+1}-\sqrt{x}
\]
при больших $x$ лучше вычислять в эквивалентной форме
\[
\frac{1}{\sqrt{x+1}+\sqrt{x}}.
\]
Порядок выполнения операций также может влиять на погрешность, например при
суммировании большого числа величин сильно различающихся масштабов.

\subsection{Баланс погрешностей}

Уменьшение шага дискретизации не всегда улучшает окончательный результат.
В численном дифференцировании типична оценка
\[
E(h)\lesssim C_1h^p+\frac{C_2\varepsilon}{h^q},
\]
где первое слагаемое --- погрешность метода, а второе отражает усиление
ошибок исходных данных и округления. Поэтому существует оптимальный
диапазон $h$; брать шаг ``как можно меньше'' неправильно.

Практическое правило состоит в согласовании уровней ошибок: обычно нет
смысла решать дискретную задачу с точностью, намного превосходящей
неопределённость исходных данных и самой математической модели.


\subsection{Прямая и обратная погрешности}

Для численного алгоритма полезно различать \textbf{прямую погрешность} ---
расстояние от вычисленного результата до точного решения --- и
\textbf{обратную погрешность}: насколько надо изменить исходные данные, чтобы
вычисленный результат стал точным решением слегка возмущённой задачи.

Для системы $Ax=b$ и приближения $\widehat x$ невязка равна
\[
    r=b-A\widehat x.
\]
Поскольку
\[
    A(x-\widehat x)=r,
\]
получаем
\[
    \|x-\widehat x\|\le \|A^{-1}\|\,\|r\|.
\]
При $b\ne0$ отсюда следует стандартная оценка
\[
    \boxed{
    \frac{\|x-\widehat x\|}{\|x\|}
    \le
    \kappa(A)\frac{\|r\|}{\|b\|}.
    }
\]
Поэтому малая невязка ещё не гарантирует малую ошибку решения, если задача плохо
обусловлена. Обратно устойчивый алгоритм стремится давать ответ, который является
точным решением близкой по исходным данным задачи.

\subsection{Алгебраическая погрешность итерационного решения}

После дискретизации часто возникает система
\[
    A_h x_h=b_h,
\]
которую решают итерационно. Тогда к погрешности дискретизации добавляется
\textbf{алгебраическая погрешность} из-за преждевременной остановки итераций.
Рациональный критерий остановки согласуют с точностью дискретизации: нет смысла
уменьшать невязку на много порядков ниже уровня ошибки сеточного метода.

Для стационарной итерации
\[
    x^{k+1}=Bx^k+c
\]
сходимость при любом $x^0$ эквивалентна условию
\[
    \rho(B)<1.
\]
Таким образом, устойчивость и скорость решения алгебраической задачи также
входят в общий бюджет погрешности вычислительной модели.

\subsection{Практическая оценка порядка и ошибки дискретизации}

Если
\[
    Q_h=Q+C h^p+o(h^p),
\]
то по трём расчётам на сетках $h$, $h/r$, $h/r^2$ можно оценить наблюдаемый
порядок
\[
    p_{\rm obs}
    \approx
    \frac{\ln\left|\dfrac{Q_h-Q_{h/r}}
    {Q_{h/r}-Q_{h/r^2}}\right|}{\ln r}.
\]
Совпадение $p_{\rm obs}$ с теоретическим порядком --- важная часть
верификации. Если порядок не проявляется, возможны слишком грубая сетка,
недостаточная гладкость решения, доминирующая погрешность граничных условий,
алгебраическая ошибка или ошибка реализации.

\subsection{Погрешность и неопределённость --- не одно и то же}

В вычислительном моделировании полезно различать:
\begin{itemize}
    \item \textbf{погрешность} --- различие между вычисленным и некоторым
    определённым точным значением;
    \item \textbf{неопределённость} --- неполное знание входных данных,
    параметров или самой структуры модели.
\end{itemize}
Например, дискретизационную ошибку можно уменьшать сгущением сетки, а
неопределённость экспериментально измеренного модуля упругости --- только за счёт
дополнительной информации о материале. В итоговом прогнозе обе составляющие
необходимо учитывать раздельно.

\subsection{Контроль законов сохранения как диагностика}

Для моделей, основанных на законах баланса, сильным диагностическим тестом
является контроль массы, импульса, энергии или другого инварианта. Нарушение
дискретного баланса может указывать на ошибку схемы, слишком большой шаг,
некорректные граничные условия или программную ошибку. Такой контроль не заменяет
оценку нормы ошибки, но часто обнаруживает проблемы, которые незаметны по одной
локальной величине.

\subsection{Что важно сказать на экзамене}

Нужно уверенно различать три основных класса: \textbf{погрешность модели и
данных, погрешность метода, вычислительную погрешность}. Отдельно следует
различать обусловленность задачи и устойчивость алгоритма. Для вычислений
с плавающей точкой важно помнить модель относительного округления и
явление потери значащих цифр.

\newpage
